% =====================================================================
%  APPENDIX: ROBUSTNESS ANALYSIS  (output-gap equation, V14 main spec)
%  Requires: booktabs. Place after \appendix.
%  All figures verified against the full re-run on 2026-06-15.
% =====================================================================
\section{Robustness Analysis}\label{app:robustness}

This appendix documents the full battery of checks behind the preferred
output-gap specification (V14, eq.~\ref{eq:output_est}): the country
fixed-effects estimator
\[
y_{ik}=\rho_y\,y_{i,k-1}+\alpha_S\,S_{ik}
       +\alpha_{\text{above}}\,F^{CP,\text{above}}_{i,k-2}
       +\alpha_{\text{below}}\,K^{CP,\text{below}}_{ik}
       +\alpha_{DI}\,F^{DI}_{i,k-1}
       +\alpha_{S\cdot DI}\,S_{ik}F^{DI}_{i,k-1}
       +\gamma_i+\varepsilon_{ik},
\]
estimated on 38 OECD countries over $t\in[4,14]$ (Q4.2019--Q2.2022,
$N=418$), with standard errors clustered by country (CRV1, small-sample
corrected). The exercises span six concerns: (i)~identification and
collinearity, (ii)~estimator and dynamic-panel bias, (iii)~heterogeneous-
treatment / negative-weight contamination, (iv)~time effects and sample
composition, (v)~instrument construction and functional form, and
(vi)~small-cluster inference. Table~\ref{tab:robustness_summary} lists each
check, what it probes, and its outcome.

The two headline results are robust throughout: the above-the-line
capacity-preservation multiplier ($\alpha_{\text{above}}\approx0.54$, flow,
two-quarter lag) and the demand-injection ``push-on-string'' slope
($\alpha_{S\cdot DI}\approx-0.041$, break-even at $S^\star\approx36$) retain
their sign, magnitude and significance across every specification, sample
window and inference method. The below-the-line stock channel
($\alpha_{\text{below}}\approx0.26$) is the only fragile coefficient: it is
sensitive to the inclusion of time trends and to the pre-pandemic
zero-stock anchor, and is therefore reported as a theoretically motivated
lower bound rather than a precisely identified estimate.

\begin{table}[htbp]
\centering
\footnotesize
\renewcommand{\arraystretch}{1.15}
\caption{Robustness battery for the V14 output-gap specification}
\label{tab:robustness_summary}
\begin{tabular}{@{}p{0.255\textwidth} p{0.37\textwidth} p{0.315\textwidth}@{}}
\toprule
\textbf{Check} & \textbf{What it probes} & \textbf{Outcome} \\
\midrule
\multicolumn{3}{@{}l}{\textit{A.\ Identification \& collinearity}}\\
Within/between decomposition & Variation usable by country FE & Within-share $>77\%$ for all regressors ($S$: 94\%, $F^{CP}$: 97\%). \\
VIF \& pairwise correlations & Collinearity among Above, Below, $S$, DI & Main terms VIF $<1.4$, $|r|<0.42$; the DI$\times S$ VIFs ($\approx14$) are mechanical interaction artefacts. \\
$S$--composition orthogonality & Fiscal mix endogenous to severity? & $r(\bar S,\bar F^{CP})\!\approx\!0.13$; composition driven by institutions/safety nets, not contemporaneous severity. \\
\addlinespace
\multicolumn{3}{@{}l}{\textit{B.\ Estimator \& dynamic-panel bias}}\\
\texttt{feols} vs.\ \texttt{plm} & Estimator invariance & Identical point estimates. \\
Nickell-bias assessment & Short-$T$ dynamic-panel bias & Bias $<0.7$pp, negligible vs.\ SEs; corroborated by Arellano--Bond GMM. \\
Mundlak / Chamberlain CRE & Fixed vs.\ random effects & Within-effects $=$ V14; joint test $\chi^2=25.7$, $p<0.001\Rightarrow$ FE required. \\
\addlinespace
\multicolumn{3}{@{}l}{\textit{C.\ Heterogeneous-treatment robustness}}\\
de Chaisemartin--D'Haultf\oe uille & Negative TWFE weights & Negative-weight shares: Above $8.7\%$, DI $2.7\%$ (benign); Below $26.8\%$ (flags its fragility). \\
\addlinespace
\multicolumn{3}{@{}l}{\textit{D.\ Time effects \& sample composition}}\\
Year-FE / linear trend / Quarter-FE & Common time variation & Above-Flow \& DI$\times S$ trend-robust; Below-Stock trend-confounded (sign flips under Year-FE). \\
Start/end restrictions ($t\!\ge\!5,\ge\!6,\le\!13$) & Pre-pandemic anchor & Above/DI$\times S$ stable; Below inflates to $5.87$ at $t\!\ge\!6$ --- the Q4.2019 zero-stock anchor is essential. \\
Estimation window (Narrow/Wide/Only-2020) & Window choice & Above $0.37$--$0.73$, DI$\times S$ $-0.04$ to $-0.05$ stable; Below significant only in baseline. \\
Outlier exclusion (TUR, IRL) & Inflation- / MNC-driven GDP & Core unchanged: Above $0.45^{*}$, DI$\times S$ $-0.035^{*}$. \\
\addlinespace
\multicolumn{3}{@{}l}{\textit{E.\ Instrument construction \& functional form}}\\
Take-up grid (loans 40/60/80\%, guar.\ 25/35/50\%) & Below-stock take-up assumptions & Above ($0.54$) \& DI$\times S$ ($-0.04$) invariant; Below $0.13$--$0.26$ (baseline largest). \\
Decay stock ($\delta=0$--$0.25$) & Stock-accumulation rule & Above \& DI$\times S$ stable; Below sign-sensitive to depreciation. \\
Lag selection (Above 0--3, DI 0--2) & Disbursement timing & Only Above lag-2 and DI lag-1 yield sign-correct, theory-consistent estimates. \\
Alternative containment (peak $S$, transit mobility) & Measurement of $S$ & Above and push-on-string survive both. \\
\addlinespace
\multicolumn{3}{@{}l}{\textit{F.\ Mechanism \& heterogeneity}}\\
Tightening vs.\ loosening asymmetry & Hysteresis / ratchet & $\Delta S^{+}=-0.125^{***}$ vs.\ $\Delta S^{-}=-0.005$ ($\approx26\times$): structural support for AR(1) persistence. \\
Sample splits ($S$, income, pre-debt, safety net) & Where channels operate & Above strongest in high-$S$/high-debt; DI push-on-string in all splits; Below only in high-$S$ \& weak-safety-net groups. \\
Regional groupings (G7, EU, advanced) & Composition & Consistent with baseline (power-limited at small $N$). \\
\addlinespace
\multicolumn{3}{@{}l}{\textit{G.\ Small-cluster inference} ($N_{\text{cl}}=38$)}\\
Wild-cluster restricted bootstrap-$t$ (CGM 2008, Rademacher) & CRV1 over-rejection & Above-Flow $p=0.008$; Below \& DI$\times S$ marginal ($p\approx0.06$); qualitative pattern unchanged. \\
Cluster-jackknife (leave-one-country-out) & Influential clusters & Above-Flow stable; no single country drives the result. \\
\addlinespace
\multicolumn{3}{@{}l}{\textit{H.\ Specification limit}}\\
Persistence-reduction test (CP$\times y_{i,k-1}$) & Does CP lower $\rho_y$? & Not separately identified in $T=11$; rejected, and addressed instead through the level channels. \\
\bottomrule
\end{tabular}
\vspace{4pt}
\parbox{\textwidth}{\scriptsize\textit{Notes:} Dependent variable: output gap $y_{ik}$ (pp of potential GDP). All
fiscal variables in pp of 2019 GDP; $S$ on a $0$--$100$ scale. All models include country fixed effects with
country-clustered standard errors. Significance: $^{*}p<0.1$, $^{**}p<0.05$, $^{***}p<0.01$. Above $=$
$F^{CP,\text{above}}_{i,k-2}$ (above-the-line flow); Below $=$ $K^{CP,\text{below}}_{ik}$ (below-the-line
accumulated stock); DI$\times S=$ the demand-injection $\times$ stringency interaction.}
\end{table}

Taken together, the battery establishes that the paper's two structural
claims---a positive, lagged capacity-preservation multiplier and a
state-dependent demand-injection effect that turns contractionary under
tight containment---do not depend on any single modelling choice, sample
window, instrument-construction assumption, or inference procedure. The
below-the-line stock channel is retained for theoretical completeness but
flagged as the empirically weakest element of the specification.
